\documentclass[11pt]{article}

\usepackage[margin=1in]{geometry}
\usepackage{amsmath,amssymb,amsthm}
\usepackage[colorlinks=true,linkcolor=blue,urlcolor=blue,citecolor=blue]{hyperref}

\newtheorem{theorem}{Theorem}
\newtheorem{question}{Research question}
\newcommand{\M}{\mathcal M}
\newcommand{\E}{\mathbb E}
\newcommand{\Pp}{\mathbb P}

\title{Bayes Marginals of Sufficient Statistics on Mean-Value Space}
\author{}
\date{}

\begin{document}
\maketitle

\section{Setup and interpretation}

Consider a regular minimal exponential family
\[
  P_{\beta}(x)
  =\frac{\exp\{-\beta^{\mathsf T}T(x)\}}{Z(\beta)},
  \qquad
  \mu(\beta)=\E_{\beta}[T(X)]\in\M.
\]
The canonical parameter $\beta$ and mean parameter $\mu$ are in one-to-one
correspondence in the interior of the parameter spaces.  We may therefore
write $P_{\mu}:=P_{\beta(\mu)}$ and put a prior density $w$ directly on the
mean-value space:
\[
  M\sim Q^w,\qquad Q^w(M\in B)=\int_B w(\mu)\,d\mu.
\]
Conditional on $M=\mu$, let $X_1,\ldots,X_n$ be i.i.d. from $P_\mu$ and set
\[
  \widehat\mu_n=\frac1n\sum_{i=1}^nT(X_i).
\]
The corresponding Bayes marginal (or prior-predictive distribution) is
\[
  P^{w,(n)}(x^n)
  =\int_{\M}\prod_{i=1}^nP_\mu(x_i)w(\mu)\,d\mu.
\]
This is not a posterior statement.  It compares the prior distribution of
the latent mean $M$ with the marginal distribution of $\widehat\mu_n$ before
the data are observed.  Since
\[
  \E[\widehat\mu_n\mid M]=M,
  \qquad
  \operatorname{Var}(\widehat\mu_n\mid M)=\frac{V(M)}n,
\]
we have the high-level picture
\[
  \widehat\mu_n=M+\text{centered noise of size }n^{-1/2}.
\]
Thus the marginal law of $\widehat\mu_n$ should approach the prior law of
$M$.

\section{The existing result and its logarithm}

Theorem~1 of Giuffrida, Garlaschelli and Gr\"unwald
(\href{https://arxiv.org/abs/2509.01064}{arXiv:2509.01064}) states, in a
finite-support exponential-family setting, that for a regular prior $w$ and
an interior compact convex set $B\subset\M$,
\begin{equation}\label{eq:existing}
  \left|
  P^{w,(n)}(\widehat\mu_n\in B)-Q^w(M\in B)
  \right|
  =O\!\left(\frac{\log n}{n}\right).
\end{equation}
Their proof enlarges and shrinks $B$ by a boundary buffer with squared radius
of order $\log(n)/n$.  Equivalently, its Euclidean width is
\[
  r_n\asymp\sqrt{\frac{\log n}{n}}.
\]
For parameters farther than $r_n$ from the boundary, a multivariate Chernoff
bound is applied.  Local quadratic equivalence between Kullback--Leibler
divergence and squared Euclidean distance gives
\[
  \Pp_{\mu}(\widehat\mu_n\text{ crosses the buffer})
  \lesssim \exp(-c n r_n^2)=n^{-c}.
\]
The logarithm therefore enters through the choice of a buffer large enough
to make this crossing probability polynomially small.

There is a point worth checking carefully: a boundary layer of width $r_n$
generally has prior mass of order $r_n$, not $r_n^2$.  Hence the displayed
sandwich argument appears, without additional cancellation, to give only
$O(\sqrt{\log(n)/n})$.  The faster rate in \eqref{eq:existing} may still be
true under suitable smoothness, but it seems to require an argument that uses
the centering of $\widehat\mu_n-M$, rather than bounding inward and outward
crossings separately.

\section{Assumptions to distinguish}

A sharper theorem should state explicitly:
\begin{itemize}
  \item whether $T(X)$ has finite support, is lattice-valued, or is continuous;
  \item minimality and regularity of the exponential family;
  \item that $B$ stays in a compact subset of the interior of $\M$, so that
        Fisher information is uniformly nonsingular;
  \item the regularity of $w$ (continuity, H\"older continuity, Lipschitz
        continuity, or $C^k$ smoothness);
  \item the geometry of $B$ and its boundary; and
  \item whether the bound is pointwise in $B,w$ or uniform over a class.
\end{itemize}
Continuity and strict positivity alone do not provide a quantitative modulus
of continuity, so they may be insufficient for a universal polynomial rate.
Boundary points of $\M$ also require separate treatment because $V(\mu)$ may
degenerate there.

\section{Possible sharp rate}

For a smooth test function $f$, conditional Taylor expansion suggests
\begin{equation}\label{eq:generator}
  \E f(\widehat\mu_n)-\E f(M)
  =\frac1{2n}\E\!\left[
    \operatorname{tr}\{V(M)\nabla^2f(M)\}
  \right]+\text{smaller terms}.
\end{equation}
The nominal $n^{-1/2}$ term vanishes because the conditional noise is
centered.  This suggests that, under appropriate smoothness, the logarithm in
\eqref{eq:existing} can be removed:
\[
  \left|P^{w,(n)}(\widehat\mu_n\in B)-Q^w(M\in B)\right|=O(n^{-1}).
\]
The rate $1/n$ is generically sharp.  For example, if
$X_i\mid M=\mu\sim\operatorname{Bernoulli}(\mu)$ and
$M\sim\operatorname{Uniform}(0,1)$, then the prior-predictive distribution of
$K_n=\sum_iX_i$ is uniform on $\{0,\ldots,n\}$.  Therefore, for an interior
interval $B=[a,b]$,
\[
  P^{w,(n)}(K_n/n\in B)
  =\frac{\#\{k:k/n\in B\}}{n+1}
  =(b-a)+O(n^{-1}),
\]
and lattice rounding generally prevents a uniform $o(n^{-1})$ result.

An $O(n^{-2})$ term is more naturally expected as the remainder in an
expansion
\[
  P^{w,(n)}(\widehat\mu_n\in B)-Q^w(M\in B)
  =\frac{A(B,w)}n+O(n^{-2}).
\]
The total error becomes $O(n^{-2})$ only in special cases where $A(B,w)=0$,
or after an explicit bias or lattice correction.  A bound
$O(\log\log(n)/n)$ would improve the published logarithm, but removing the
logarithm entirely is the more natural target.

\section{Possible tools and next questions}

Stein's method may provide uniform conditional normal approximations for
\[
  \sqrt n\,(\widehat\mu_n-\mu)\mid M=\mu
  \approx N(0,V(\mu)).
\]
However, an off-the-shelf central limit bound may be too coarse: the desired
$1/n$ behavior arises from cancellation after mixing over $w$, while
$1_B$ is nonsmooth at the boundary.  A useful Stein approach would preserve
the first-order cancellation, perhaps through an operator or generator
expansion such as \eqref{eq:generator}, followed by a careful smoothing or
boundary argument.

Other plausible tools are local central limit theorems, Edgeworth expansions,
saddlepoint and Laplace approximations, small-diffusion expansions, and
explicit continuity corrections for lattice families.  A related stronger
goal is a local expansion of the induced
prior-predictive density,
\[
  g_n(t)=w(t)+\frac1{2n}
  \sum_{i,j}\partial_{ij}\{V_{ij}(t)w(t)\}+o(n^{-1}),
\]
which would also yield expansions for ratios of Bayes marginals and hence
connect directly to e-values and reverse information projections.

\begin{question}
Under minimal regularity conditions on the exponential family, prior $w$ and
set $B$, what is the sharp rate and leading correction in
\[
  P^{w,(n)}(\widehat\mu_n\in B)-Q^w(M\in B)?
\]
Can the $\log n$ factor be removed, can the coefficient $A(B,w)$ be
identified, and when does it vanish?
\end{question}

This gives a compact starting point for the brainstorm: first verify the
existing rate and assumptions, then study a sharp $1/n$ theorem, its
second-order correction, boundary and lattice effects, and finally the
consequences for Bayes-marginal likelihood ratios and e-values.

\end{document}
